% GENERATED FILE. Edit canonical lessons, then run npm run generate:books.
% canonical-source-hash: fnv1a64:ae58a42f15be1761
% canonical-lessons: ZH-C426-dachuangfang, ZH-C426-keman, ZH-C426-ludian, ZH-C426-xingji, ZH-C426-chazuo

\chapter{Booking a Room}
\label{ch:zh-booking-a-room}
\hlchaptermodality{426}

\begin{chapteropening}
\textbf{By the end of this chapter:} \emph{I can ask for a standard room, hear that a hotel is full, find an inn, ask about stars and pay the bill.}

Complete the last lesson of chapter 426: I can ask for a standard room, hear that a hotel is full, find an inn, ask about stars and pay the bill.
\end{chapteropening}

\section[dàchuángfáng]{dàchuángfáng — room with a double bed}

\label{lesson:ZH-C426-dachuangfang}



\begin{quote}
Before the new one: say the Mandarin for a room key card, then the Mandarin for the hotel lobby.
\end{quote}

\subsection*{You'll want to know: dàchuángfáng}
\emph{dàchuángfáng} — \textquotedblleft{}room with a double bed\textquotedblright{}.

\begin{quote}
\emph{wǒmen yào yì jiān dàchuángfáng} — We'd like a double room.
\end{quote}

Its characters come later; for now it is a word for the ear.

\begin{cousinweb}
\emph{dàchuángfáng} is built from \zh{大} 'big', \emph{chuáng} 'a bed' and \zh{房} 'a room', all inherited from Middle Chinese, \emph{daH}, \emph{dzrjang} and \emph{bjang} (Guangyun, in Baxter's transcription).
\end{cousinweb}

\begin{grammarlens}[title={What you've built}]
One more word for getting around.
\end{grammarlens}

\subsection*{Your turn}
\begin{itemize}
\raggedright
  \item \emph{Say it:} \emph{dàchuángfáng}
  \item \emph{Say it:} \emph{dàchuángfáng}, once more
  \item \emph{Say it:} \emph{wǒmen yào yì jiān dàchuángfáng} — We'd like a double room.
  \item \emph{Recall:} say the Mandarin for a room key card, then the Mandarin for the hotel lobby, then say \emph{dàchuángfáng} again
\end{itemize}

\begin{tcolorbox}[breakable,colback=teal!4,colframe=teal!35!black,title={Before you move on}]
What does dàchuángfáng mean? (\textquotedblleft{}Room with a double bed\textquotedblright{}.) Say it once more.
\end{tcolorbox}
\section[kèmǎn]{kèmǎn — fully booked; with no vacancies}

\label{lesson:ZH-C426-keman}



\begin{quote}
Before the new one: say the Mandarin for the hotel lobby, then the Mandarin for a room with a double bed.
\end{quote}

\subsection*{You'll want to know: kèmǎn}
\emph{kèmǎn} — \textquotedblleft{}fully booked; with no vacancies\textquotedblright{}.

\begin{quote}
\emph{duìbuqǐ jīntiān kèmǎn le} — Sorry, we're full tonight.
\end{quote}

Its characters come later; for now it is a word for the ear.

\begin{cousinweb}
\emph{kèmǎn} is built from \zh{客} 'a guest' and \emph{mǎn} 'full', both inherited from Middle Chinese, \emph{khaek} and \emph{manX} (Guangyun, in Baxter's transcription), which glosses the second 'full, filled'.
\end{cousinweb}

\begin{grammarlens}[title={What you've built}]
One more word for getting around.
\end{grammarlens}

\subsection*{Your turn}
\begin{itemize}
\raggedright
  \item \emph{Say it:} \emph{kèmǎn}
  \item \emph{Say it:} \emph{kèmǎn}, once more
  \item \emph{Say it:} \emph{duìbuqǐ jīntiān kèmǎn le} — Sorry, we're full tonight.
  \item \emph{Recall:} say the Mandarin for the hotel lobby, then the Mandarin for a room with a double bed, then say \emph{kèmǎn} again
\end{itemize}

\begin{tcolorbox}[breakable,colback=teal!4,colframe=teal!35!black,title={Before you move on}]
What does kèmǎn mean? (\textquotedblleft{}Fully booked; with no vacancies\textquotedblright{}.) Say it once more.
\end{tcolorbox}
\section[\texorpdfstring{\zh{旅店} (lǚdiàn)}{旅店 (lǚdiàn)}]{\zh{旅店} (lǚdiàn) — inn; small hotel}

\label{lesson:ZH-C426-ludian}



\begin{quote}
Before the new one: say the Mandarin for a room with a double bed, then the Mandarin for fully booked.
\end{quote}

\subsection*{You'll want to know: \zh{旅店}}
\textbf{\zh{旅店}} — \emph{lǚdiàn} — \textquotedblleft{}inn; small hotel\textquotedblright{}.

\begin{quote}
\textbf{\zh{这家旅店很小}} \emph{zhè jiā lǚdiàn hěn xiǎo} — This inn is very small.
\end{quote}

\begin{cousinweb}
\textbf{\zh{旅店}} (\emph{lǚdiàn}) is built from \zh{旅} 'to travel' and \zh{店} 'an inn', both inherited from Middle Chinese, \emph{ljoX} and \emph{temH} (Guangyun, in Baxter's transcription); the dictionary glosses the first as a body of troops, so the sense 'to travel' is not given there.
\end{cousinweb}

\begin{grammarlens}[title={What you've built}]
One more word for getting around.
\end{grammarlens}

\subsection*{Your turn}
\begin{itemize}
\raggedright
  \item \emph{Say it:} \emph{lǚdiàn}
  \item \emph{Say it:} \emph{lǚdiàn}, once more
  \item \emph{Say it:} \emph{zhè jiā lǚdiàn hěn xiǎo} — This inn is very small.
  \item \emph{Recall:} say the Mandarin for a room with a double bed, then the Mandarin for fully booked, then say \emph{lǚdiàn} again
\end{itemize}

\begin{tcolorbox}[breakable,colback=teal!4,colframe=teal!35!black,title={Before you move on}]
What does \zh{旅店} mean? (\textquotedblleft{}Inn; small hotel\textquotedblright{}.) Say it once more.
\end{tcolorbox}
\section[xīngjí]{xīngjí — hotel star rating}

\label{lesson:ZH-C426-xingji}



\begin{quote}
Before the new one: say the Mandarin for fully booked, then the Mandarin for a small inn.
\end{quote}

\subsection*{You'll want to know: xīngjí}
\emph{xīngjí} — \textquotedblleft{}hotel star rating\textquotedblright{}.

\begin{quote}
\emph{zhè jiā jiǔdiàn de xīngjí hěn gāo} — This hotel has a lot of stars.
\end{quote}

Its characters come later; for now it is a word for the ear.

\begin{cousinweb}
\emph{xīngjí} is built from \zh{星} 'a star' and \emph{jí} 'a rank', both inherited from Middle Chinese, \emph{seng} and \emph{kip} (Guangyun, in Baxter's transcription).
\end{cousinweb}

\begin{grammarlens}[title={What you've built}]
One more word for getting around.
\end{grammarlens}

\subsection*{Your turn}
\begin{itemize}
\raggedright
  \item \emph{Say it:} \emph{xīngjí}
  \item \emph{Say it:} \emph{xīngjí}, once more
  \item \emph{Say it:} \emph{zhè jiā jiǔdiàn de xīngjí hěn gāo} — This hotel has a lot of stars.
  \item \emph{Recall:} say the Mandarin for fully booked, then the Mandarin for a small inn, then say \emph{xīngjí} again
\end{itemize}

\begin{tcolorbox}[breakable,colback=teal!4,colframe=teal!35!black,title={Before you move on}]
What does xīngjí mean? (\textquotedblleft{}Hotel star rating\textquotedblright{}.) Say it once more.
\end{tcolorbox}
\section[chāzuò]{chāzuò — power socket}

\label{lesson:ZH-C426-chazuo}



\begin{quote}
Before the new one: say the Mandarin for a small inn, then the Mandarin for the hotel star rating.
\end{quote}

\subsection*{You'll want to know: chāzuò}
\emph{chāzuò} — \textquotedblleft{}power socket\textquotedblright{}.

\begin{quote}
\emph{fángjiān lǐ de chāzuò zài nǎr} — Where's the socket in the room?
\end{quote}

Its characters come later; for now it is a word for the ear.

\begin{cousinweb}
\emph{chāzuò} is built from \emph{chā} 'to insert' and \emph{zuò} 'a seat', both inherited from Middle Chinese, \emph{tsrheap} and \emph{dzwaH} (Guangyun, in Baxter's transcription), which glosses the second as the seat of a bed.
\end{cousinweb}

\begin{grammarlens}[title={What you've built}]
One more word for getting around.
\end{grammarlens}

\subsection*{Your turn}
\begin{itemize}
\raggedright
  \item \emph{Say it:} \emph{chāzuò}
  \item \emph{Say it:} \emph{chāzuò}, once more
  \item \emph{Say it:} \emph{fángjiān lǐ de chāzuò zài nǎr} — Where's the socket in the room?
  \item \emph{Recall:} say the Mandarin for a small inn, then the Mandarin for the hotel star rating, then say \emph{chāzuò} again
\end{itemize}

\begin{tcolorbox}[breakable,colback=teal!4,colframe=teal!35!black,title={Before you move on}]
What does chāzuò mean? (\textquotedblleft{}Power socket\textquotedblright{}.) Say it once more.
\end{tcolorbox}
